\chapter{The Fatou and Julia Sets}
\label{chp:the_fatou_and_julia_sets}

The informal discussion in Chapter~\ref{chp:examples} suggested that the
dynamics of a rational map $R$ induces a subdivision of the complex sphere into
two sets, namely the Fatou set $F$ and the Julia set $J$ of $R$. Here, we
introduce the notion of equicontinuity of a family of maps, and then formally
define the Fatou and Julia sets in terms of the equicontinuity of the family
$\{R^n\}$. In addition, we consider conditions which imply that a family of
analytic maps is equicontinuous (or a normal family) in a domain.

% --- SECTION ------------------------------------------------------------------
% --- THE FATOU AND JULIA SETS -------------------------------------------------
\section{The Fatou and Julia Sets}
\label{sec:the_fatou_and_julia_sets}

\begin{definition}[Continuous function]
    Given two metric spaces $(X, d_X)$ and $(Y, d_Y)$. A map $f: X \to Y$ is
    \emph{continuous at $x_0$} in $X$ if and only if, for every positive
    $\varepsilon$ there exists a positive $\delta$ such that for all $x$ in $X$,
    \begin{equation*}
        d_X(x_0, x) < \delta \quad \text{implies} \quad d_Y(f(x_0), f(x)) < \varepsilon.
    \end{equation*}
    The function $f$ is \emph{continuous on a subset $X_0$ of $X$} if it is
    continuous at each point $x_0$ of $X_0$.
\end{definition}

Clearly, in this definition, $\delta$ depends on $f$, $x_0$ and $\varepsilon$.
For equicontinuity we will look for a $\delta$ that works for all $f$ in a
family of functions.

\begin{definition}[Equicontinuous family]
    Given two metric spaces $(X, d_X)$ and $(Y, d_Y)$. A family $\mathcal{F}$ of
    maps of $(X, d_X)$ into $(Y, d_Y)$ is \emph{equicontinuous at $x_0$} in $X$
    if and only if, for every positive $\varepsilon$ there exists a positive
    $\delta$ such that for all $x$ in $X$, and for all $f$ in $\mathcal{F}$,
    \begin{equation*}
        d_X(x_0, x) < \delta \quad \text{implies} \quad d_Y(f(x_0), f(x)) < \varepsilon.
    \end{equation*}
    The family $\mathcal{F}$ is \emph{equicontinuous on a subset $X_0$ of $X$}
    if it is equicontinuous at each point $x_0$ of $X_0$.
\end{definition}

It is important to realize that the equicontinuity of $\mathcal{F}$ at $x_0$ is
the formal expression for the idea of preservation of proximity which was
introduced informally in Chapter~\ref{chp:examples}: indeed, it implies that
every function in $\mathcal{F}$ maps the open ball with center $x_0$ and radius
$\delta$ into a ball of radius at most $\varepsilon$.

Suppose $\mathcal{F} = \{R^n \mid n \in \N\}$ is an equicontinuous family, then
given $z \in \ExtC$, points $w$ close to $z$ will remain close after $n$
iterations ($R^n(w)$ will be close to $R^n(z)$). Indeed, if we take
$\varepsilon_n$ to be the maximum distance we want between the $n$-th iterates,
we can obtain a maximum distance between the $(n-1)$-th iterates
($\varepsilon_{n-1}$). Doing this recursively we can obtain a maximum distance
$\varepsilon_0 = \delta$ between $w$ and $z$ so that their $n$-th iterates are
as close as we want, see Figure~\ref{fig:equicontinuity_for_iterates}.
\begin{figure}
    \centering
    \resizebox{\textwidth}{!}{\begin{tikzpicture}[
    >-Stealth,
    point/.style={circle, fill=black, inner sep=1pt},
    dashed circle/.style={dashed, draw=black, minimum size=3.5cm, circle},
    mapping arrow/.style={->, thick, shorten >=5pt, shorten <=5pt}
]

    % --- REGION 1: N ----------------------------------------------------------
    \node[dashed circle] (C1) {};
    \node[below right=1cm and 1.25cm of C1.center] {$< \delta$};
    \node[above right=0.5cm and -1.2cm of C1.center] {$N$};
    
    % Draw Blob
    \pic [scale=0.35] at (C1.center) {blob7};
    
    % Points inside N
    \node[point, label=above:{$z$}] (z1) at ($(C1.center) + (0, 0)$) {};
    \node[point, label=left:{$w$}] (w1) at ($(C1.center) + (0.1, -0.7)$) {};

    % --- REGION 2: R(N) -------------------------------------------------------
    \node[dashed circle, right=2cm of C1] (C2) {};
    \node[below right=1cm and 1.25cm of C2.center] {$< \varepsilon_1$};
    \node[above right=0.85cm and -0.85cm of C2.center] {$R(N)$};

    % Draw Blob
    \pic [scale=0.25, rotate=45] at (C2.center) {blob4};
    
    % Points inside R(N)
    \node[point, label=above:{$R(z)$}] (z2) at ($(C2.center) + (0, 0)$) {};
    \node[point, label=above:{$R(w)$}] (w2) at ($(C2.center) + (-0.65, -0.5)$) {};


    % --- REGION 3: R^n(N) -----------------------------------------------------
    \node[dashed circle, right=5.4cm of C2] (C3) {};
    \node[below right=1cm and 1.25cm of C3.center] {$< \varepsilon_n$};
    \node[above right=-0.1cm and 0.35cm of C3.center] {$R^n(N)$};
    
    % Draw Blob
    \pic [scale=0.24, rotate=-80] at (C3.center) {blob3};
    
    % Points inside R^n(N)
    \node[point, label=below:{$R^n(z)$}] (z3) at ($(C3.center) + (0, 0)$) {};
    \node[point, label=above:{$R^n(w)$}] (w3) at ($(C3.center) + (-0.35, 0.4)$) {};


    % --- ARROWS & LABELS ------------------------------------------------------
    \draw[mapping arrow] (C1.east) -- (C2.west) node[midway, above=3pt] {\Large $R$};
    
    \coordinate (Arrow2End) at ($(C2.east) + (1.9, 0)$);
    \draw[mapping arrow] (C2.east) -- (Arrow2End) node[midway, above=3pt] {\Large $R$};
    
    \node[right=0.4cm of Arrow2End, font=\Huge, inner sep=0pt] (dots) {$\dots$};
    \draw[mapping arrow] ($(dots.east) + (0.2, 0)$) -- (C3.west) node[midway, above=3pt] {\Large $R$};
\end{tikzpicture}}
    \caption{Equicontinuity for the family $\mathcal{F} = \{R^n\}_{n \in \N}$.}
    \label{fig:equicontinuity_for_iterates}
\end{figure}

\begin{remark}[Equicontinuity is preserved under union]
    If the family $\mathcal{F}$ is equicontinuous on each of the subsets
    $D_\alpha$ of $X$, then it is automatically equicontinuous on the union
    $\cup D_\alpha$.
\end{remark}

Taking the collection $\{D_\alpha\}$ to be the class of all open subsets of $X$
on which $\mathcal{F}$ is equicontinuous, this leads to the following general
principle.

\begin{theorem}[$\exists$ a maximal open subset where $\mathcal{F}$ is equicontinuous]
    Let $\mathcal{F}$ be any family of maps, each mapping $(X, d_X)$ into $(Y,
    d_Y)$. Then there is a maximal open subset of $X$ on which $\mathcal{F}$ is
    equicontinuous.
\end{theorem}

\begin{corollary}[$\exists$ a maximal open subset where $\{f^n\}$ is equicontinuous]
    If $f$ maps a metric space $(X, d_X)$ into itself, then there is a maximal
    open subset of $X$ on which the family of iterates $\{f^n \mid n \in \N\}$
    is equicontinuous.
\end{corollary}

This (at last) provides us with a formal definition for the Fatou and Julia sets
of a rational map $R$, and the following terminology is chosen to honor the
creators of the theory.

\begin{definition}[Fatou and Julia sets]
    Let $R$ be a non-constant rational function. The \emph{Fatou set} of $R$ is
    the maximal open subset of $\ExtC$ on which $\{R^n \mid n \in \N\}$ is
    equicontinuous, and the \emph{Julia set} of $R$ is its complement in
    $\ExtC$.

    We denote the Fatou set of a rational map $R$ by $F$, and the Julia set by
    $J$, although sometimes, when we need to mention $R$ explicitly, we use
    $F(R)$ and $J(R)$ instead.
\end{definition}

\begin{remark}[$F$ is open and $J$ is compact]
    Note that, by definition, $F(R)$ is \emph{open}, and $J(R)$ is
    \emph{compact}.
\end{remark}

Although the term \enquote{Julia set} is standard, the use of \enquote{Fatou
set} was suggested as late as 1984 (in \cite{Blanchard1984}). It seems
appropriate, but the reader should be familiar with the common alternatives,
namely \emph{the stable set}, and \emph{the set of normality}. This latter term
is a reference to normal families of analytic functions, but we have preferred
to base our definition on equicontinuity because of its more immediate geometric
appeal.

We end this section with two simple, but important, properties of $F$ and $J$. A
rational map, and in particular, a Möbius map and its inverse, satisfy a
Lipschitz condition with respect to the spherical (and chordal) metric on
$\ExtC$ (Theorem~\ref{thm:rational_maps_satify_a_lipschitz_condition}). Thus, if
we conjugate $R$ with respect to a Möbius map, equicontinuity is transferred in
the obvious way, and we have

\begin{theorem}[Equicontinuity is preserved under conjugation]
    Let $R$ be a non-constant rational map, let $g$ be a Möbius map, and let $S
    = g R g^{-1}$. Then $F(S) = g(F(R))$ and $J(S) = g(J(R))$.
\end{theorem}

A similar argument leads to

\begin{theorem}[$F$ and $J$ are preserved under iteration]
    \label{thm:F_and_J_are_preserved_under_iteration}
    For any non-constant rational map $R$, and any positive integer $p$, $F(R^p)
    = F(R)$ and $J(R^p) = J(R)$.
\end{theorem}
\begin{proof}
    See \cite[51]{Beardon1991}.
\end{proof}

% --- SECTION ------------------------------------------------------------------
% --- COMPLETELY INVARIANT SETS ------------------------------------------------
\section{Completely invariant sets}
\label{sec:completely_invariant_sets}

\begin{definition}[Invariant set under $g$]
    Given a map $g$ of a set $X$ into itself, a subset $E$ of $X$ is:
    \begin{enumerate}[label=(\alph*)]
        \item \emph{forward invariant} if $g(E) = E$;
        \item \emph{backward invariant} if $g^{-1}(E) = E$;
        \item \emph{completely invariant} if $g(E) = E = g^{-1}(E)$.
    \end{enumerate}
\end{definition}

\begin{remark}[$g$ surjective $\implies$ backward invariant $=$ completely invariant]
    If $g$ is surjective (that is, if $g(X) = X$) then the concepts of backward
    invariance and complete invariance coincide. Indeed, by surjectivity (but
    not necessarily in the general case, because there may be some points in $E$
    where $g^{-1}(x) = \emptyset$),
    \begin{equation*}
        g(g^{-1}(E)) = E
    \end{equation*}
    which, with backward invariance, yields $g(E) = E$. The requirement of
    surjectivity here is crucial and without it, there is a difference between
    backward and complete invariance.
\end{remark}

\begin{example}
    The map $z \mapsto \exp(z)$ shows how these concepts can be different. Under
    this map $\C$ is backward invariant, but not completely invariant.
\end{example}

As our main concern is rational maps, we recall that these map $\ExtC$ onto
itself: thus for rational maps, backward invariance coincides with complete
invariance.

We illustrate these ideas with the following result which will be of importance
later.

\begin{theorem}[Completely invariant under $R$ and finite $\implies$ at most two elements]
    \label{thm:completely_invariant_implies_non_empty_non_singleton}
    Let $R$ be a rational map of degree at least two and suppose that a finite
    set $E$ is completely invariant under $R$. Then $E$ has at most two
    elements.
\end{theorem}
\begin{proof}
    See \cite[52]{Beardon1991}.
\end{proof}

It is convenient to collect together several, general, but elementary,
properties of completely invariant sets and throughout, \emph{we shall assume
that $g: X \to X$ is surjective} so that we may use complete invariance and
backward invariance interchangeably.

\begin{proposition}[Complete invariance is preserved under conjugation]
    Given a surjective map $g$ of a set $X$ into itself, and a completely
    invariant subset $E$ of $X$ under $g$. Then, if $h$ is a bijection of $X$
    onto itself, then $h(E)$ is completely invariant under $h g h^{-1}$.
\end{proposition}

\begin{proposition}[Complete invariance is preserved under intersection]
    Given a surjective map $g$ of a set $X$ into itself, and a family of
    completely invariant subsets $\{E_\alpha\}$ of $X$ under $g$. Then,
    \begin{equation*}
        E = \bigcap E_\alpha
    \end{equation*}
    is completely invariant.
\end{proposition}
\begin{proof}
    The proof can be obtained by applying the standard construction which is
    used to generate, for example, subgroups, subspaces and topologies. The
    operator $g^{-1}$ commutes with the intersection operator, that is, for any
    collection $\{E_\alpha\}$ of sets,
    \begin{equation*}
        g^{-1}\left(\bigcap E_\alpha \right) = \bigcap g^{-1}(E_\alpha)
    \end{equation*}
\end{proof}

This means, that we can take any subset $E_0$ and form the intersection, say
$E$, of all those completely invariant sets which contain $E_0$: obviously, $E$
is then the \emph{smallest completely invariant set that contains $E_0$}.

\begin{definition}[Generation of a completely invariant set]
    We say that $E_0$ \emph{generates} $E$ if $E$ is the intersection of all the
    completely invariant sets that contain $E_0$. That is, $E$ is then the
    smallest completely invariant set that contains $E_0$.
\end{definition}

We now introduce an equivalence relation which greatly facilitates the
discussion of completely invariant sets.

\begin{definition}[Relation $\sim$ (under $g$)]
    Given a map $g$ of a set $X$ into itself, for any $x$ and $y$ in $X$, we
    define the relation $\sim$ on $X$ by $x \sim y$ if and only if there exist
    non-negative integers $n$ and $m$ with
    \begin{equation}
        \label{eq:equivalence_relation}
        g^n(x) = g^m(y).
    \end{equation}
\end{definition}

\begin{remark}[$\sim$ is an equivalence relation]
    Obviously the relation $\sim$ is symmetric and reflexive, and it is also
    transitive, since \eqref{eq:equivalence_relation} and $g^p(y) = g^q(z)$
    imply that
    \begin{equation*}
        g^{n+p}(x) = g^{m+q}(z);
    \end{equation*}
    thus $\sim$ is an equivalence relation on $X$.
\end{remark}

\begin{definition}[Orbit of $x$ (under $g$)]
    We denote the equivalence class containing $x$ by $[x]$, and we call this
    the \emph{orbit} of $x$ (some authors call this the \enquote{great orbit}).
\end{definition}

It is easy to identify the orbit $[x]$, and we have

\begin{theorem}[{$[x]$ is the completely invariant set generated by $\{x\}$}]
    Let $g: X \to X$ be any map of $X$ onto itself. Then $[x]$ is the completely
    invariant set generated by the singleton $\{x\}$.
\end{theorem}
\begin{proof}
    See \cite[53]{Beardon1991}.
\end{proof}

\begin{corollary}[Completely invariant $\iff$ union of orbits]
    A set $E$ is completely invariant if and only if it is a union of
    equivalence classes $[x]$.
\end{corollary}

This suggests we should also look at the interior, closure and boundary
operators.

\begin{theorem}[Complete invariance is preserved under the complement, interior and boundary operators]
    \label{thm:complete_invariance_is_preserved_under_complement_interior_and_boundary}
    Let $g$ be any continuous, open map of a topological space $X$ onto itself
    and suppose that $E$ is completely invariant. Then so are the complement $X
    - E$, the interior $\interior{E}$, the boundary $\boundary{E}$ and the
    closure $\closure{E}$, of $E$.
\end{theorem}
\begin{proof}
    See \cite[53]{Beardon1991}.
\end{proof}

We now prove that the division of the sphere into Fatou and Julia sets is
completely invariant: this is fundamental.

\begin{theorem}[$F$ and $J$ are completely invariant]
    \label{thm:F_and_J_are_completely_invariant}
    Let $R$ be any rational map. Then the Fatou set $F$ and the Julia set $J$
    are completely invariant under $R$.
\end{theorem}
\begin{proof}
    See \cite[54]{Beardon1991}.
\end{proof}

We can say more than Theorem~\ref{thm:F_and_J_are_completely_invariant} when $R$
is a polynomial, and we have

\begin{theorem}[$F_\infty$ is completely invariant for polynomials]
    \label{thm:F_infty_is_completely_invariant_for_polynomials}
    Let $P$ be a polynomial of degree at least two. Then $\infty$ is in $F(P)$,
    and the component $F_\infty$ of $F$ containing $\infty$ is completely
    invariant under $P$.
\end{theorem}
\begin{proof}
    See \cite[54]{Beardon1991}.
\end{proof}

We end this section with a simple observation which will be used several times
later.

\begin{proposition}[Finite number of components $\implies$ they are periodic (under a continuous map $f: X \to X$)]
    \label{prp:finite_number_of_Fatou_components_implies_they_are_periodic}
    Let $f$ be a continuous map of a topological space $X$ onto itself, and
    suppose that $X$ has only a finite number of components $X_j$. Then for some
    integer $m$, each $X_j$ is completely invariant under $f^m$.
\end{proposition}
\begin{proof}
    See \cite[55]{Beardon1991}.
\end{proof}

This, of course, was the essential idea in the proof of
Theorem~\ref{thm:completely_invariant_implies_non_empty_non_singleton}.

% --- SECTION ------------------------------------------------------------------
% --- NORMAL FAMILIES AND EQUICONTINUITY ---------------------------------------
\section{Normal families and equicontinuity}
\label{sec:normal_families_and_equicontinuity}

We have defined the Fatou and Julia sets in terms of the equicontinuity of the
family $\{R^n\}$. However, equicontinuity is closely related to normal families
of analytic functions, and once we have understood this connection (which is
expressed in the Arzelà-Ascoli Theorem), we can exploit the more powerful
results about normal families of analytic functions to derive further
information about the Fatou and Julia sets. This is our program for this
section.

We begin with the definitions relevant to normal families.

\begin{definition}[Pointwise convergence]
    A sequence $(f_n)$ of maps from a metric space $(X, d_X)$ to a metric space
    $(Y, d_Y)$ \emph{converges pointwise on $X$} to some map $f$ if for each $x$
    in $X$: for every positive $\varepsilon$ there exists a natural number $N$
    such that for all $n \geq N$,
    \begin{equation*}
        d_Y(f_n(x), f(x)) < \varepsilon.
    \end{equation*} 
\end{definition}

\begin{definition}[Uniform convergence]
    A sequence $(f_n)$ of maps from a metric space $(X, d_X)$ to a metric space
    $(Y, d_Y)$ \emph{converges uniformly on $X$} to some map $f$ if for every
    positive $\varepsilon$ there exists a natural number $N$ such that for all
    $x$ in $X$, and for all $n \geq N$,
    \begin{equation*}
        d_Y(f_n(x), f(x)) < \varepsilon.
    \end{equation*}
\end{definition}

\begin{remark}[Uniform convergence $\implies$ pointwise convergence]
    Note that for a family to be uniformly convergent is a stronger statement
    than for it to be pointwise convergent. Indeed, if we write the definitions
    in terms of limits we have:
    \begin{alignat*}{2}
        &\lim_{n \to \infty} f_n = f \text{ pointwise} &&\iff \\
        &&&\forall x \in X \lim_{n \to \infty} d_Y(f_n(x), f(x)) = 0; \\
        &\lim_{n \to \infty} f_n = f \text{ uniformly} &&\iff \\
        &&&\lim_{n \to \infty} \sup\{d_Y(f_n(x), f(x)) \mid \forall x \in X\} =
        0.
    \end{alignat*}
    This last expression may also clarify why in
    section~\ref{sec:a_topology_on_the_rational_functions}, we called $\rho$ the
    \emph{metric of uniform convergence}. Since, using that metric, convergence
    in $\mathcal{C}$ is the same as uniform convergence.
\end{remark}

\begin{definition}[Local uniform convergence]
    A sequence $(f_n)$ of maps from a metric space $(X, d_X)$ to a metric space
    $(Y, d_Y)$ \emph{converges locally uniformly on $X$} to some map $f$ if each
    point $x$ of $X$, has a neighborhood on which $f_n$ converges uniformly to
    $f$.
\end{definition}

\begin{remark}[Local uniform convergence $\implies$ uniform convergence in compacts]
    Given a sequence $(f_n)$ of maps from a uniform metric space $(X, d_X)$ to a
    metric space $(Y, d_Y)$ that converges locally uniformly on $X$. Then, it
    converges uniformly on each compact subset of $X$.
\end{remark}

\begin{definition}[Normal family]
    A family $\mathcal{F}$ of maps from a metric space $(X, d_X)$ to a metric
    space $(Y, d_Y)$ is said to be \emph{normal}, or a \emph{normal family}, in
    $X$ if every infinite sequence of functions from $\mathcal{F}$ contains a
    subsequence which converges locally uniformly on $X$.
\end{definition}

We give the usual reminder that the limit of the convergent subsequence need not
lie in $\mathcal{F}$. Although we have formulated these definitions for general
metric spaces, we shall only be concerned with subdomains of the complex sphere
with the chordal (or spherical) metric, and we come now to a statement of the
Arzelà-Ascoli Theorem.

\begin{theorem}[In $\ExtC$ with continuous functions: equicontinuous family $\iff$ normal family]
    Let $D$ be a subdomain of the complex sphere, and let $\mathcal{F}$ be a
    family of continuous maps of $D$ into the sphere. Then $\mathcal{F}$ is
    equicontinuous in $D$ if and only if it is a normal family in $D$.
\end{theorem}
\begin{proof}
    See \cite[222]{Ahlfors1979}.
\end{proof}

As an illustration of the direct use of normality rather than equicontinuity, we
consider Vitali's Theorem: this exploits analytic continuation and so enables us
to derive information about the limit $f$ of a sequence $(f_n)$ throughout a
domain $D$ simply from a knowledge of $f$ on only a small part of $D$.

\begin{theorem}[Vitali's Theorem]
    \label{thm:Vitalis_theorem}
    Suppose that the family $\{f_1, f_2, \ldots\}$ of analytic maps is normal in
    a domain $D$, and that $f_n$ converges pointwise to some function $f$ on
    some non-empty open subset $W$ of $D$. Then $f$ extends to a function $F$
    which is analytic on $D$, and $f_n \to F$ locally uniformly on $D$.
\end{theorem}
\begin{proof}
    See \cite[56]{Beardon1991}.
\end{proof}

We turn now to obtain conditions which guarantee the normality of a family of
analytic functions.

\begin{proposition}[In $\ExtC$ with continuous functions: holomorphic and uniformly bounded $\iff$ equicontinuous family $\iff$ normal family]
    Given a family of functions $\mathcal{F}$ which are holomorphic and
    uniformly bounded in $D$. Then, $\mathcal{F}$ is equicontinuous, and hence
    normal in $D$.
\end{proposition}
\begin{proof}
    It is an elementary consequence of Cauchy's Integral Formula that if $f$ is
    holomorphic and satisfies $|f| \leq M$ in a domain $D$ in $\C$, then $|f'|$
    is bounded above on each compact subset $K$ of $D$, the upper bound
    depending on $M$, $K$, and $D$ \emph{but not on $f$}.
\end{proof}

It is possible to obtain a variety of mild generalizations of this result
without much effort, but all such results are hopelessly inadequate for our
purposes. Instead, we pass these by and go straight to one of the most powerful
results concerning normality: this result is absolutely essential for the study
of iteration of rational maps, yet its substantial proof can safely be omitted
without prejudicing an understanding of iteration theory.

\begin{theorem}[Montel's fundamental normality criterion]
    \label{thm:analytic_maps_not_touching_zero_one_and_infty_form_a_normal_family}
    Let $D$ be a domain on the complex sphere $\ExtC$, and let $\Omega$ be the
    domain $\ExtC - \{0, 1, \infty\}$. Then the family $\mathcal{F}$ of all
    analytic maps $f: D \to \Omega$ is normal in $D$.
\end{theorem}
\begin{proof}
    See \cite[59]{Beardon1991}.
\end{proof}

The hypothesis of
Theorem~\ref{thm:analytic_maps_not_touching_zero_one_and_infty_form_a_normal_family}
is that each function in $\mathcal{F}$ fails to take the values $0$, $1$ and
$\infty$ at any point of $D$. Let's see how
Theorem~\ref{thm:analytic_maps_not_touching_zero_one_and_infty_form_a_normal_family}
can be used to produce two variations on the same theme. In the first one of
these, we assume only that each function $f$ fails to take three values which
now \emph{may depend on $f$}. It is too much to expect that three values can be
completely arbitrary, and we must add the requirement that they are uniformly
separated on the sphere.

\begin{theorem}[Montel's theorem with separated values]
    \label{thm:analytic_maps_not_touching_af_bf_and_cf_form_a_normal_family}
    Let $\mathcal{F}$ be a family of maps, each analytic in a domain $D$ on the
    complex sphere. Suppose also that there is a positive constant $m$ and, for
    each $f$ in $\mathcal{F}$, three distinct points $a_f$, $b_f$ and $c_f$ in
    $\ExtC$ such that:
    \begin{enumerate}[label=(\roman*)]
        \item $f$ in $\mathcal{F}$ does not take the values $a_f$, $b_f$ and
        $c_f$ in $D$; and
        \item $\min\{\sigma(a_f, b_f), \sigma(b_f, c_f), \sigma(c_f, a_f)\} \geq
        m$.
    \end{enumerate}
    Then $\mathcal{F}$ is normal in $D$.
\end{theorem}
\begin{proof}
    See \cite[58]{Beardon1991}.
\end{proof}

Of course,
Theorem~\ref{thm:analytic_maps_not_touching_af_bf_and_cf_form_a_normal_family}
contains
Theorem~\ref{thm:analytic_maps_not_touching_zero_one_and_infty_form_a_normal_family}
as a special case, for we can take $a_f = 0$, $b_f = 1$ and $c_f = \infty$ for
every $f$. There is another variation of
Theorem~\ref{thm:analytic_maps_not_touching_zero_one_and_infty_form_a_normal_family}
that we shall need.

\begin{theorem}[Montel's theorem for omitted functions]
    Let $D$ be a domain, and suppose that the functions $\varphi_1$, $\varphi_2$
    and $\varphi_3$ are analytic in $D$, and are such that the closures of the
    domains $\varphi_j(D)$ are mutually disjoint. Now let $\mathcal{F}$ be a
    family of functions, each analytic in $D$, such that for every $z$ in $D$,
    and every $f$ in $\mathcal{F}$, $f(z) \neq \varphi_j(z)$, $j = 1$, $2$, $3$.
    Then $\mathcal{F}$ is normal in $D$.
\end{theorem}
\begin{proof}
    See \cite[58]{Beardon1991}.
\end{proof}